\chapter{Conclusões}

Este trabalho teve como objetivo realizar uma modelagem matemática da rede de transporte público por ônibus de São Paulo, operada pela SPTrans, buscando analisar a eficiência da malha e identificar pontos de melhoria. A análise abrangeu um período temporal estratégico, de 10 anos, contemplando os impactos da pandemia de Covid-19 na demanda, e procurou preencher as lacunas de estudos anteriores ao focar na integração de múltiplas fontes de dados e na estrutura geoespacial detalhada do sistema.

A metodologia empregada, centrada na teoria dos grafos e no uso de ferramentas computacionais para ingestão, tratamento e visualização dos dados, provou-se útil para o tratamento, padronização e unificação de um volume alto de dados. Obtivemos sucesso na integração dos dados estruturais do GTFS e das séries temporais de passageiros transportados. Conseguimos também juntar dados demográficos complementares do Censo do IBGE de 2022, enfrentando, porém, dificuldades computacionais para integrar totalmente os dados do censo com as informações geográficas do GTFS. Ainda assim, enfrentamos alguns desafios de inconsistências e dados faltantes, frequentemente encontrados em fontes públicas, ao solicitar o histórico completo à SPTrans, e conseguimos, com sucesso, documentar algumas inconsistências que ainda restaram nos dados que estudamos.

Alcançamos os objetivos propostos com a criação do grafo da rede municipal, que viabilizou análises de centralidade e nos permitiu observar a melhoria que ocorreu no sistema de transporte ao longo do tempo. Também conseguimos analisar rotas com alto volume de passageiros, e elaborar uma análise temporal das principais estatísticas das bases obtidas, na qual observamos mudanças significativas na estrutura do sistema de transporte paulista e da demanda dos usuários. Adicionalmente, o estudo de caso sobre a reestruturação das linhas da Zona Oeste, originado da inauguração da Estação Vila Sônia, demonstrou a importância dos dados para avaliar os impactos das modificações na infraestrutura urbana, fornecendo argumentos para possíveis novos estudos e políticas. 

Ao avaliar a eficiência da rede de transporte público de São Paulo, entendida neste trabalho como a capacidade de equilibrar oferta e demanda, sem priorizar o lucro das operadoras de transporte, é possível perceber um cenário de contrastes: de um lado, há avanços importantes, como o aumento do uso do Bilhete Único, a ampliação das gratuidades e a descentralização do sistema, que mostram um progresso tecnológico, social e estrutural. Porém, ao analisarmos casos como a reestruturação das linhas na Zona Oeste, depois da abertura da Estação Vila Sônia, conseguimos identificar um aumento de custos do transporte para o usuário final, além da observação de desequilíbrio entre oferta e demanda, como na linha 745810 (Jd. Boa Vista), onde a quantidade de viagens é claramente insuficiente. Além disso, identificamos no Jardim Ângela e em Itaquera dois exemplos de regiões que necessitam de mais planejamento e foco em políticas públicas, de forma a melhorar a qualidade do transporte. Essas situações mostram a necessidade de revisar o planejamento e as políticas, para garantir um serviço mais equilibrado e eficiente para toda a cidade, que melhore a qualidade de vida dos moradores de São Paulo, sempre se amparando em dados, como os discutidos neste trabalho.

A principal contribuição prática desta pesquisa é a elaboração de uma base de dados confiável, que viabiliza análises mais aprofundadas do sistema, permitindo, por exemplo, a criação de um painel interativo, que fornece ao cidadão uma visão mais aprofundada de um dos tópicos que mais impacta o seu dia a dia: o transporte. 

Esse painel centraliza as bases de dados tratadas e as análises realizadas, permitindo que servidores públicos e cidadãos explorem visualmente informações operacionais de qualquer linha, como trajeto, frequência programada e histórico de demanda de passageiros. O painel fornece transparência e subsídio técnico, possibilitando a população fundamentação técnica para solicitações de melhoria no sistema de transporte.

Desse modo, essa monografia avança na compreensão da mobilidade urbana paulistana ao fornecer não apenas uma análise estrutural e temporal da rede de ônibus, mas também uma base de dados e plataforma dinâmica para o planejamento e aperfeiçoamento contínuo do transporte público, visando um sistema mais eficiente e condizente com as necessidades da população.